\section{Introduction}
\label{sec:intro}

A common safety story for chat models is that they \emph{refuse}: asked to do something harmful, they decline. Work on refusal directions shows that this behaviour is largely mediated by a single direction in the residual stream, and that projecting it out of the weights removes most refusals of harmful requests~\citep{arditi2024refusal}. Agents face a different situation. Nobody asks them to cause harm. Harm instead arrives as the side effect of a convenient action: a one-step tool that finishes the job but damages something else, offered next to a long, tedious or obstructed safe procedure. Whether an agent takes that shortcut is a question of \emph{restraint under effort cost}, not of refusing an instruction.

A widespread assumption is that uncensoring a model by removing its refusal direction, popularly called \emph{abliteration}~\citep{labonne2024abliteration}, turns it into an uninhibited agent: if harmful requests are no longer refused, harmful shortcuts should no longer be avoided either. This paper asks whether refusal and restraint are the same mechanism. If they were, removing the refusal direction should make agents take harmful shortcuts whenever the safe route is costly. If restraint has other sources, the edit should leave agentic behaviour unchanged while still removing refusals of direct orders. Our data fit neither pole. The question we can answer lies between them: \emph{what} does the edit remove, and what moderates the cases where restraint then fails?

We edit Qwen3.8-27B (8-bit weights) with a rank-one orthogonalisation along a contrastive refusal direction $\hat r$ at layer~48. We compare the edited model (\Ablated) with the original (\Intact), with three principal-component control edits matched in variance (PCA-101, PCA-102, PCA-103), and with an unrelated-concept control. Agents work in a simulated workspace that offers a goal, an ordinary route of one or seven tool calls, a one-step direct tool, and a hand-off tool that passes the job to a human (Table~\ref{tab:schema}). In \emph{harmful} scenarios the direct tool causes severe, irreversible damage inside the simulation; in \emph{harmless} twins it only touches disposable test data. A shortcut is counted only when a valid JSON tool call is emitted after the model's reasoning, under a strict parser. Nothing is inferred from prose.

\paragraph{Thesis.} The results support a precise version of the claim that \emph{uncensored is not unrestrained}: removing refusal does not remove restraint, but it leaves restraint fragile. They organise into three takeaways.
\begin{enumerate}[leftmargin=1.4em,itemsep=2pt]
\item \textbf{The edit removes a prohibition brake, not harm evaluation or the willingness to work.} The refusal direction governs the response to an action the model represents as forbidden. Ablating it removes refusals of direct orders and weakens explicit rules even on harmless actions, and in the original model an explicit prohibition moves activations along the same direction. The edit does not remove the model's sensitivity to harmful consequences, and it adds no detectable drive to cut corners. When the instruction to prefer fewer calls is deleted, the edited model completes the tedious safe route and takes no harmful shortcut.
\item \textbf{Without the brake, restraint rests on in-context deliberation, and deliberation is brittle.} Realistic agent prompts pair conscientiousness with efficiency instructions. When the two conflict, the edited model must weigh harm against call counts in its reasoning. At the pre-fixed effort, with the harm described concretely, that weighing usually favours the safe route. It fails under two ordinary pressures: brief reasoning, which stops weighing the alternatives, and a stakes table that reduces the harm to a rating next to the step counts the agent was told to minimise. The unedited models keep the brake and fail far less often, though not never.
\item \textbf{A decision made during reasoning is hard to read before it and easy to pass on through it.} With thinking on, pre-thinking activations add nothing detectable beyond the task cell; with thinking off, where the model must commit at once, they carry a small within-cell signal. Fine-tuning on the edited teacher's transcripts transfers its shortcut tendency even when every demonstrated action is safe, and paraphrasing the teacher's reasoning removes the transfer.
\end{enumerate}
We call reasoning effort, stakes format, the efficiency instruction and distillation \emph{moderators}, not mediators: we manipulate them and observe the outcome, but we do not identify the internal pathway through which each acts.

\paragraph{Contributions.} The evidence for these takeaways comes from nine experiment blocks and two deconfounder slices (summarised in Table~\ref{tab:overview}).
\begin{itemize}
\item We qualify a refusal-ablated checkpoint family against a direct-instruction positive control, and we document why a second candidate family failed (Section~\ref{sec:phase1}).
\item \emph{First takeaway.} We show a dissociation between action spaces. On direct orders the edit removes refusal, yet the edited model still escalates whenever a hand-off tool is listed. In autonomous tasks agents walk the long safe route rather than escape, and removing the hand-off tool changes nothing. The edited model's few harmful shortcuts at the pre-fixed effort depend on the fewer-calls instruction and do not reproduce item by item. The edit also weakens explicit rules on harmless actions, and the refusal direction responds to explicit prohibition in the original model (Sections~\ref{sec:direct} and~\ref{sec:probes}).
\item \emph{Second takeaway.} We identify the pressures under which restraint fails. Lower reasoning effort reduces direct-order compliance but triples autonomous harmful shortcuts. An abstract stakes table produces the largest failures, in every arm, and a deconfounder shows this is an interaction of tabular layout with the loss of the concrete description (Section~\ref{sec:direct}).
\item \emph{Third takeaway.} We find that pre-thinking layer-48 probes add no detectable information beyond the task cell when the model reasons, and a small increment when it does not (Section~\ref{sec:probes}). The shortcut tendency transfers through distillation on safe-route demonstrations, but not through paraphrased or unrelated reasoning (Section~\ref{sec:distill}).
\item We characterise reviewer contagion, two-agent coordination and monitor oversight. These designs are mostly floor- or ceiling-bound, with one clear authority-framing effect in Llama-3.3-70B (Section~\ref{sec:multi}).
\end{itemize}

\subsection{Related work}
\paragraph{Refusal directions and activation steering.} Representation engineering reads and controls high-level behaviours through directions in activation space~\citep{zou2023repe}, and activation addition steers generation without optimisation~\citep{turner2023actadd,rimsky2024caa}. \citet{arditi2024refusal} showed that refusal is mediated by a single direction whose ablation removes refusals while preserving capability. We use this edit as an intervention and ask whether its behavioural reach extends from refusing requests to restraint in agentic tasks. Later work finds that refusal is better described by multi-dimensional concept cones~\citep{wollschlager2025geometry}, and that models encode harmfulness separately from refusal~\citep{zhao2025harmfulness}. That separation is one reason an edit to the refusal direction need not remove harm-sensitive behaviour. Jailbreak research documents complementary failure modes of safety training~\citep{wei2023jailbroken}.

\paragraph{Agentic safety and control.} Avoiding negative side effects of goal pursuit is a long-standing safety problem~\citep{amodei2016concrete,krakovna2020sideeffects}; our harmful shortcut is such a side effect. Refusal-vector ablation applied to Llama-3.1-70B agents yields agents that perform harmful tasks~\citep{lermen2024agents}, and refusal-trained models are easily jailbroken once deployed as browser agents~\citep{kumar2024browser}. Both results concern agents asked to do harm; we ask about harm as a by-product of a benign goal. LM-emulated sandboxes~\citep{ruan2024toolemu} and harmful-agent benchmarks~\citep{andriushchenko2024agentharm} measure whether agents carry out harmful tasks. Text-game benchmarks measure trade-offs between reward and ethical behaviour~\citep{pan2023machiavelli}, and in-context feedback loops can drive reward hacking~\citep{pan2024feedback}. Frontier models can also pursue goals through covert actions~\citep{meinke2024scheming}. The AI-control agenda studies monitoring protocols that remain safe even against an untrusted model~\citep{greenblatt2024control}. Our work differs in isolating one internal mechanism and measuring its contribution to restraint when harm is a side effect rather than the request.

\paragraph{Reasoning, faithfulness and monitoring.} Chain-of-thought explanations can be unfaithful to the factors that drive a model's answer~\citep{turpin2023unfaithful,lanham2023faithfulness}, and monitoring reasoning can be undermined by optimisation pressure~\citep{baker2025monitoring}. Longer test-time reasoning can also degrade behaviour~\citep{gema2025inverse}. Our reasoning-effort results show that the sign of the effect can differ between direct orders and autonomous tasks. Linear probes are our read-out tool~\citep{alain2016probes}.

\paragraph{Distillation and trait transfer.} Knowledge distillation transfers behaviour from teacher to student~\citep{hinton2015distilling}, and LoRA makes such fine-tuning cheap~\citep{hu2021lora}. \citet{cloud2025subliminal} report that behavioural traits can pass through data that appears unrelated to them. Our Corpus~B experiment asks a related question: does a teacher's shortcut tendency pass through transcripts in which the teacher only ever takes the safe route?

\paragraph{Methodological commitments.} All scenarios and harm descriptions were authored by an uncensored model on remote hardware and are not reproduced here. Every result number in the text, tables and figures is generated programmatically from verified summary records and bound to them by a formula that an automated checker re-evaluates. We use exact paired McNemar tests~\citep{mcnemar1947}, Wilson intervals~\citep{wilson1927}, Newcombe intervals for paired differences~\citep{newcombe1998}, Holm correction~\citep{holm1979}, cluster-level sign-flip tests and Firth-penalised logistic models~\citep{firth1993}. The total compute cost was \$\fact{budget.billed} of a \$\fact{budget.cap} budget. Section~\ref{sec:limits} collects every threat to validity we know of.
